% Research cards: BAS, PER
% Checked on the web on September 23, 2026, with two 2025 BDGD (V11) actually opened.
% "(to be confirmed)" marks what was not verified in a source or in data.

\section{Network spatial base and non-technical losses}

\subsection*{Inside the BDGD}

Two 2025 BDGD (version V11) were downloaded and opened for this research
[2]. The Instruction Manual could not be downloaded (the item on the portal
responds with an error), so the semantics of some fields were inferred from the
order of magnitude of the values, and this is indicated.

\paragraph{Access.} ANEEL Open Data, ODbL license, annual update
[1]. One geodatabase per distribution utility on the ANEEL ArcGIS Hub, listable
through a search API and downloadable \textbf{without authentication} [2]. 2025
sizes: from about 1~MB (cooperatives) to 4~GB (Cemig-D). The ODbL has a
share-alike clause for derived databases used publicly:
validate with legal counsel whether the product layers fall under it (to be confirmed).

\paragraph{Lag.} In a 2025 BDGD, the period runs from January to
December 2025, with extraction in May and publication in August 2026: in that
edition, publication occurred \textbf{8 to 20 months} after the reference
month. The public data serve to
build the network and the pilot, not for monthly operation.

\begin{table}[H]
\centering\small
\caption{BDGD entities relevant to the catalog.}
\begin{tabularx}{\textwidth}{P{3.2cm} L}
\toprule
\textbf{Entity} & \textbf{What it carries (verified in the files)} \\
\midrule
CTMT (feeder) & substation, voltages, HV transformer; \textbf{monthly circuit energy} (ENE\_01 to 12, compatible with injected energy, to be confirmed); \textbf{technical losses} by segment and transformation (PERD\_*); \textbf{monthly non-technical losses} at medium and low voltage (PNTMT, PNTBT) \\
UNTRMT (transformer) & rated power, iron and total losses, feeder, substation, electrical set, municipality, phases, connection date \\
UCBT, UCMT (units) & no geometry of their own; link to the point (PONNOT) and to the transformer with 100\% fill rate in the sample; class, CNAE, load-curve typology, installed load, DG code, \textbf{monthly energy, DIC and FIC}; MV also with demand. \textbf{The unit code is a 64-character hash} \\
BE, EP, PT, PNT & energy balance, pass-through energy, monthly technical and non-technical losses of the distribution utility \\
CONJ, ARAT, SUB & polygons of electrical set, service area and substation (the latter came empty in one of the files) \\
SSDMT, SSDBT & network segments, with connection points to build the graph \\
\bottomrule
\end{tabularx}
\end{table}

\paragraph{Pitfalls verified in the data.}
\begin{itemize}
  \item \textbf{Different units across distribution utilities}: feeder
    energy appears in kWh in one and in MWh in the other, and within a single
    table there are fields on different scales. Every loss/energy ratio must
    be normalized by distribution utility and by field, or the naive calculation gives
    ``technical loss of 250\%''.
  \item \textbf{The BDGD non-technical loss appears to be apportionment}: round values
    that match the aggregate table indicate accounting allocation, not
    measurement (to be confirmed in the Manual).
  \item \textbf{Implausible MMGD}: in a distribution utility in SC, annual energy
    over installed power reaches 4,470~kWh/kW/year, incompatible with
    photovoltaic generation (median 900), which is compatible with the
    injection and unit problems.
  \item \textbf{Codes}: feeder, transformer and connection points
    appear to be operational codes, but the consumer unit is anonymized.
    Linking client data \emph{by unit} requires the distribution utility to
    provide the mapping; linking \emph{through the network} tends to work (to be
    confirmed with a client).
\end{itemize}

\paragraph{ANEEL loss indicators [5--8].} In 2024, total loss was
about 14\% of injected energy: technical 7.4\% (44.6~TWh) and non-technical
6.6\% (40.2~TWh). The ten largest distribution utilities account for 74\% of the
non-technical loss; Light and Amazonas Energia, 34\%. The cost of the actual non-technical loss was
estimated by ANEEL at R\$~10.3 billion. The regulatory limit comes from PRORET Submódulo
2.6 through benchmarking of socioeconomic complexity [9]. \textbf{Since 2025,
non-technical loss has been calculated on the measured market, no longer the
billed one} (Despacho 1.220/2025), which breaks the historical series. The regulatory
technical loss is calculated by ANEEL by building the BDGD in OpenDSS [12].

% ---------------------------------------------------------------------------
\subsection{BAS \textperiodcentered{} Network spatial base}

\begin{ficha}{BAS-01 \textperiodcentered{} Georeferenced electrical hierarchy}{MVP}
\campo{Data} BDGD: substation, feeder (CTMT), transformer (UNTRMT),
units (UCBT, UCMT) and their points (PONNOT); segments (SSDMT, SSDBT) for
connectivity. Converters to OpenDSS: \emph{bdgd-tools} and
\emph{bdgd2dss} [11].
\campo{Approaches} Baseline: ETL with GDAL to GeoParquet or PostGIS.
Recommended: graph from the connection points with validation of radiality and
reachability from the feeder head; unit normalization
by distribution utility; composite key (distribution utility, code); versioning by
base year.
\campo{Metrics} Percentage of units linked to transformer, point and
feeder; reachable segments; loops per circuit; divergence between
years.
\campo{Pitfalls} Substation polygon may come empty; files of several
GB; resubmissions with different dates; the BDGD is a simplified model of the network
[1].
\end{ficha}

\begin{ficha}{BAS-02 \textperiodcentered{} Administrative and regulatory boundaries}{MVP}
\campo{Data} IBGE, 2022 tract mesh and aggregates [13]; the municipality in the
BDGD uses the 7-digit IBGE code; electrical sets in the CONJ polygon of the
BDGD and in the continuity indicators [14].
\campo{Approaches} Bridge table unit, tract, municipality, set,
computed once by spatial join.
\campo{Pitfalls} 2010 and 2022 tracts incompatible; sets change
between years; matching the BDGD set code with that of the indicators (to be
confirmed).
\end{ficha}

\begin{ficha}{BAS-03 \textperiodcentered{} Data quality and coverage by entity}{MVP}
\campo{Typical BDGD problems} Disconnected sections, inconsistent
buses, wrong links, incorrect phase, inadequate impedances,
misallocated load; a single parameterization error changed the calculated technical
loss by up to 1.45\% [12].
\campo{Approaches} Completeness and validity rules by field (Great
Expectations, pandera) and a \textbf{score by entity} with physical rules
(total loss greater than iron loss; yield below the cap; loss
smaller than energy), consistency rules (sum of the units smaller than the feeder
energy), unit detection by order of magnitude and missing links.
The score becomes a layer and qualifies the others.
\campo{Pitfalls} Rules calibrated on one distribution utility fail on another because
of the units; the quality layer can be read as an accusation against the
client.
\end{ficha}

\begin{ficha}{BAS-04 \textperiodcentered{} Topology and phase correction}{V3}
\campo{Data} Voltage series from smart meters (level 3) and the BDGD.
Portaria Normativa MME 126/2026 mandates additional deployment in 2\% of
units per year for 24 months starting in March 2026, prioritizing reduction
of non-technical loss [17].
\campo{Approaches} Baseline: correlation of the voltage series.
Recommended: clustering on voltage correlation with a transformer capacity
constraint, generating candidates for field validation
(Short 2013 [19]; Hoogsteyn et al.\ 2022: voltage series outperform
power series [20]).
\campo{Pitfalls} Sparse coverage per transformer; synchronization between
meters; MMGD alters the voltage profile.
\end{ficha}

\begin{ficha}{BAS-05 \textperiodcentered{} Registry plausibility audit}{V3}
\campo{Data} BDGD generating units (power, energy, DG code);
ANEEL MMGD registry [15]; NASA POWER [16].
\campo{Approaches} Baseline: annual yield cap by municipality
from irradiation. Recommended: expected monthly range from a physical model
(pvlib), compared with the recorded injection, which is a lower bound on generation;
cross-check the BDGD power against the registry to detect W and kW. Real example:
4,470~kWh/kW/year in SC.
\campo{Pitfalls} Injection depends on self-consumption: the test works as a cap,
not as a floor; remote compensation; missing DG code.
\end{ficha}

% ---------------------------------------------------------------------------
\subsection{PER \textperiodcentered{} Non-technical losses}

\begin{ficha}{PER-01 \textperiodcentered{} Estimated loss by transformer and feeder}{MVP}
\campo{Data} Public: feeder energy, technical and non-technical losses
from CTMT, energy by unit, iron and total losses of the transformers.
Client: monthly billing (the BDGD is annual and lagged) and transformer
metering when available.
\campo{Approaches} \textbf{Baseline by feeder, with public data
only}: feeder energy minus the sum of the units (LV, MV,
public lighting) minus injected generation, compared with the BDGD
losses. By transformer without metering: technical loss from iron plus
copper proportional to the square of the load; the non-technical loss is left as the residual.
Recommended: power flow in OpenDSS for the technical loss (as ANEEL does)
and balance for the total, with a Bayesian hierarchical model that splits the
feeder residual among the transformers using the attributes from
PER-02. Recent literature: technical loss without power flow from
sparse voltage [41]; estimation of rates by region in a Brazilian city
[43].
\campo{Metrics} Error against metered transformers; stability between
cycles; agreement with the regulatory total.
\campo{Pitfalls} Mixed units; BDGD non-technical loss appears to be apportionment;
billing cycle different from the calendar month; estimated public lighting;
energy injected by DG enters the balance.
\end{ficha}

\begin{ficha}{PER-02 \textperiodcentered{} Irregularity probability}{MVP}
\campo{Data} Client: inspections with outcome (positive and negative),
billing, meter replacements, disconnections. BDGD: class, CNAE, installed load,
DIC and FIC, DG code. Context: census tract, transformer loss
(PER-01), neighborhood.
\campo{Approaches} Baseline: rules and gradient boosting on the
inspections. Recommended: boosting with \textbf{neighborhood variables} [23],
selection bias correction by inverse inspection propensity weighting,
calibration on the exploratory fraction and PU learning.
\campo{What the literature says} Glauner et al.\ showed covariate shift in a
Brazilian dataset of 3.6 million clients and 820 thousand inspections: inspections
concentrated in neighborhoods and classes make the models biased [22,24].
Coma-Puig and Carmona propose measuring recovered energy and explaining [25,26].
Revelas, Boldea and Werker (2025) show that selecting only the most probable
leads to inconsistent learning and that randomized selection allows
updating the model as if the sample were random [35]: this is the basis
adopted for the exploratory fraction. In PU learning [31--33], the assumption that the
labeled positives were chosen at random \textbf{does not hold} in losses,
because inspections are targeted: use variants with propensity dependent
on the variables (SAR-PU) [34].
\campo{Brazilian cases} Light: satellite and BDGD for illegal connections
[37], ANEEL R\&D with GESEL/UFRJ on areas with severe operating
restrictions [39]; Optimum-Path Forest (UNESP, Unicamp) [28]; survey of
28 experts [36]; irrigators in RS [44].
\campo{Metrics} Precision at top-k with $k$ equal to the inspection capacity;
energy-weighted recall; calibration \textbf{measured on the exploratory
fraction}; AUC only as a diagnostic.
\campo{Pitfalls} Metrics inflated by oversampling and leakage; inspection
negative is noisy; drift after meter replacement; LGPD.
\end{ficha}

\begin{ficha}{PER-03 \textperiodcentered{} Estimated recoverable energy}{MVP}
\campo{Approaches} Baseline: probability × estimated diverted energy
(consumption before the drop or median of the peers). Recommended: two stages,
probability × severity, with severity regressed on the confirmed
positives, optimizing net return (recovered minus inspection
cost), as in Massaferro, Di Martino and Fernández (UTE, Uruguay) [27].
\campo{Pitfalls} Severity is observed only in the positives (another bias);
accounting recovery is not sustained recovery; double counting with the
transformer balance.
\end{ficha}

\begin{ficha}{PER-04 \textperiodcentered{} Inspection plan with routes}{MVP}
\campo{Approaches} Baseline: top by expected value and nearest-neighbor
route. Recommended: \textbf{team orienteering problem} (routing with
prizes) maximizing expected value under each team's work shift, solved
with OR-Tools; \textbf{exploratory fraction drawn at random by stratum with
recorded probabilities} (5 to 15\%) [35], used later as weights in
PER-02, PER-05 and PER-06. No specific literature on
routing for losses was found.
\campo{Metrics} Value captured over the optimum; inspections per team and day;
cost per productive inspection; percentage of the plan executed.
\campo{Pitfalls} Safety criteria take precedence over the random draw;
teams that do not follow the random draw; grouping targets by route reintroduces geographic bias.
\end{ficha}

\begin{ficha}{PER-05 \textperiodcentered{} Method evaluation by region (A/B)}{V2}
\campo{Approaches} Randomization \textbf{by cluster} (transformer, neighborhood
or team) to avoid contamination; random arm as the prevalence
reference; Horvitz--Thompson estimators; effect measured in recovered
energy and in feeder loss over time (difference in
differences).
\campo{Pitfalls} Low statistical power with low prevalence; deterrence
effect among neighbors; change of the regulatory metric in 2025.
\end{ficha}

\begin{ficha}{PER-06 \textperiodcentered{} Geographic and socioeconomic bias monitor}{V2}
\campo{Approaches} Compare, by breakdown, the flagging rate with the
confirmation rate \textbf{estimated on the random fraction}; calibration and true
positive rate by group; proportion tests with correction for
multiple comparisons; Bayesian smoothing in small tracts. No
specific work on fairness in losses in Brazil was found.
Regulation recognizes socioeconomic complexity as a determinant, so
different prevalences are expected [6].
\campo{Pitfalls} Without a random arm, bias cannot be separated from real
prevalence; re-identification in fine breakdowns; using income as a variable and auditing
the same variable.
\end{ficha}

\subsection*{References for this section}
{\small
\begin{enumerate}[label={[\arabic*]},leftmargin=2.2em,itemsep=1pt]
\item ANEEL, BDGD: \url{https://dadosabertos.aneel.gov.br/dataset/base-de-dados-geografica-da-distribuidora-bdgd}
\item ANEEL ArcGIS Hub, search API: \url{https://dadosabertos-aneel.opendata.arcgis.com/}
\item BDGD Instruction Manual (not downloaded)
\item ANEEL, PRODIST Módulos 7 and 10
\item ANEEL, Perdas de Energia: \url{https://www.gov.br/aneel/pt-br/assuntos/distribuicao/perdas-de-energia}
\item ANEEL/STR, Perdas de Energia Elétrica na Distribuição (2025, base year 2024)
\item ANEEL, Luz na Tarifa dashboard, losses: \url{https://portalrelatorios.aneel.gov.br/luznatarifa/perdasenergias}
\item ANEEL, tariffs and economic and financial information
\item ANEEL, PRORET Submódulo 2.6
\item ANEEL, losses report (2026)
\item bdgd-tools: \url{https://github.com/eniocc/bdgd-tools}; bdgd2dss: \url{https://github.com/ArthurGS97/bdgd2dss}
\item Norven, BDGD quality and technical loss calculation
\item IBGE, census tract mesh
\item ANEEL, collective continuity indicators
\item ANEEL, MMGD installations
\item NASA POWER, daily API: \url{https://power.larc.nasa.gov/docs/services/api/temporal/daily/}
\item Portaria Normativa MME 126/2026
\item Enel SP, smart meters (to be confirmed)
\item Short (2013), \emph{IEEE Trans.\ Smart Grid} 4(2):651--658
\item Hoogsteyn et al.\ (2022): \url{https://arxiv.org/abs/2204.06372}
\item Glauner et al.\ (2017), survey, \emph{IJCIS}: \url{https://arxiv.org/abs/1606.00626}
\item Glauner et al.\ (2017), ISAP: \url{https://arxiv.org/abs/1702.03767}
\item Glauner et al.\ (2016), neighborhood, BDCAT: \url{https://arxiv.org/abs/1607.00872}
\item Glauner, Valtchev and State (2018), ESANN
\item Coma-Puig and Carmona (2019), \emph{Energies} 12(9):1748
\item Coma-Puig and Carmona (2022), \emph{Machine Learning}, doi:10.1007/s10994-021-06051-1
\item Massaferro, Di Martino and Fernández, \emph{IEEE TPWRS}, doi:10.1109/TPWRS.2019.2928276
\item Ramos, Souza, Papa and Falcão, OPF; Passos Júnior et al.\ (2016), \emph{EPSR} 140
\item Messinis and Hatziargyriou (2018), \emph{EPSR} 158
\item Zheng et al.\ (2018), \emph{IEEE TII} 14(4)
\item Elkan and Noto (2008), KDD
\item Kiryo et al.\ (2017), nnPU: \url{https://arxiv.org/abs/1703.00593}
\item Bekker and Davis (2020), \emph{Machine Learning} 109: \url{https://arxiv.org/abs/1811.04820}
\item SAR-PU (KU Leuven): \url{https://github.com/ML-KULeuven/SAR-PU}
\item Revelas, Boldea and Werker (2025): \url{https://arxiv.org/abs/2509.18739}
\item Savian et al.\ (2022), \emph{IJEEP} 12(3)
\item Gremes et al.\ (2024), NTL-Unet, \emph{Sensors} 24:4924
\item Assignment of buildings to units, SEGAN (2024)
\item GESEL/UFRJ and Light, ANEEL R\&D
\item Cenário Energia, Light (August 2026)
\item Bariya and Flaspohler (2025): \url{https://arxiv.org/abs/2506.21311}
\item Nguyen et al.\ (2023), \emph{Big Data}
\item Non-technical loss rates by region, \emph{EPSR} (2023)
\item Losses among irrigators, \emph{Energies} 16:6832 (2023)
\end{enumerate}}
